\documentclass[portuguese,msc,fctuc]{uc-thesis}

\import{}{acronyms.tex}

\renewcommand\nomgroup[1]{%
 % \itemsep10pt % Adjust the value as desired for the spacing between items
 \vspace{\baselineskip} % Adjust the value as desired for the spacing before a group
 \item[\large\bfseries 
  \ifstrequal{#1}{Z}{Acrónimos}{% % não é necessário devido ao pacote de acrónimos
  \ifstrequal{#1}{A}{}{%
  \ifstrequal{#1}{G}{Símbolos Gregos}{%
  \ifstrequal{#1}{S}{Sobrescritos}{% 
  \ifstrequal{#1}{U}{Subscritos}{%   
  \ifstrequal{#1}{X}{Outros Símbolos}{}}}}}}]%
 }
 
% Reducing spacing before the first group in the nomenclature to adjust for the spacing in the others
\renewcommand{\nompreamble}{\vspace{-10pt}} % Adjust the value as desired for the spacing

\makenomenclature

\begin{document}
%% Use Roman numerals for the page numbers of the title pages and table of contents.
\frontmatter
\renewcommand{\chaptermark}[1]{\markboth{\thechapter.\ #1}{}}

\import{}{thesis_info.tex}
%\import{}{nomenclature.tex}

\makecover

\maketitlepage

\makededication

\makeacknowledgement

\makeabstract

\tableofcontents

\cleardoublepage

\listoffigures

\cleardoublepage

\listoftables

%\cleardoublepage

\listofacronyms

\import{}{nomenclature}
\printnomenclature

\mainmatter
\renewcommand{\chaptermark}[1]{\markboth{\thechapter.\ #1}{}}

%% Add chapters to be imported
%\fontsize{12}{12}\selectfont
\import{chapters/}{chapter-1_Intro.tex}
\import{chapters/}{chapter-2_Fundamentos.tex}
\import{chapters/}{chapter-5_Conclusoes.tex}


\bibliography{thesis.bib}


\appendix
%% Add appendixes to be imported
\import{chapters/}{appendixA.tex}

%\annex
%% Add annexes to be imported
%\import{chapters/}{annexA.tex}

\end{document}